\section{Introduction}

Birds are a key part of the ecosystems around the world. They have important roles
that help keep the balance in nature. Different bird species have various roles in the
ecosystem such as predators, pollinators, scavengers, and seed dispersers. All these roles
are vital and have direct or indirect effects on humans~\cite{mariyappan2023ecological}.
In addition, birds are also reliable bioindicators of ecosystem health. Changes in their
distribution, population, or behaviour may be early signs of environmental
degradation~\cite{anderle2024efficiency}. Birds can play a critical role in predicting or
helping identify ecosystems facing problems such as climate change or pollution. Tracking
them can act as an early warning system.

Decline in bird population disrupts food chains and weakens the ecosystem. As humans
depend on nature for basic needs such as food production, it is crucial to track and protect
the environment~\cite{ellissoto2025biologging}. This is where the need for a feasible way
to track organisms appears. Manual methods have been used in the identification of bird
species, including visual assessments, field notes, and analyses by trained ornithologists.
However, manual methods are subjective, time-consuming, and prone to
error~\cite{kambali2024automated}.

Citizen science is a relatively recent concept in the scientific domain, but it has the
power to generate vast amounts of high-quality data~\cite{baker2021verification}. It is
the concept of letting non-professional volunteers participate in scientific data
collection~\cite{fraisl2022citizen, encarnacao2021citizen}. With digital platforms and
mobile technologies advancing, most people have access to upload their sightings. This
approach allows researchers to gather large-scale data from around the world. Gathering
this data would be expensive and impractical using traditional fieldwork. Citizen science
allows continuous tracking of organisms across regions.

Participation of volunteers also promotes public awareness as well as contributing to
scientific research~\cite{bonney2015citizen}. People engage with nature around them and
have a chance to learn about their surroundings. Large datasets constructed using citizen
science help researchers around the world.

There are existing studies that train and compare different machine learning models 
for bird species classification \cite{bilgin2022bird,tang2025comparative} however the developed models are rarely
used in practical real-world applications. This project aims to address this gap by developing 
a complete citizen science platform that integrates the trained machine learning model.

In this project, we aim to develop a citizen science platform where users can observe
birds around them, learn about the species they observe, and upload their observations.
The platform is aimed towards volunteers and enthusiasts about birdwatching. People can
interact with each other and socialise, allowing a community of like-minded people to be
formed. The platform will support citizen science and create an opportunity to monitor
bird populations. There are some fundamental challenges of citizen science such as data
quality, observer bias, and validation~\cite{johnston2022outstanding}. The platform aims
to address these challenges.

Another aim of this project is to conduct a comparative analysis of different machine
learning models for image classification. The developed citizen science platform will have
a machine learning model integrated into it to predict the observed bird species from input
images. This is an important aspect for labelling the data being generated by the public.
Creating datasets for classifying species is a lengthy process. The use of deep learning and
machine learning techniques in this process provides significant time savings and higher
accuracy. This project will compare different models and select the one that is most
suitable.

To achieve this aim, this project will conduct a comparison of deep learning architectures
using a consistent dataset. The model performance will be evaluated using metrics such as
accuracy, precision, recall, F1-score, and inference speed. This will help identify the
optimal model that will be used in the final platform.

Visual classification of bird species is considered a complex Fine-Grained Visual
Classification (FGVC) problem~\cite{chen2024fetfgvc}. This process stems from
phenotypic variations within a single species and distinguishing characteristics between
related species. Organisms can look drastically different even if they are the same species
due to factors such as age and sex. Convolutional Neural Networks (CNNs), a type of
machine learning model, are required for processing visual patterns such as colour, shape,
and texture~\cite{lolge2025evolution}. This is necessary for the accurate classification of
these patterns. The main objective of this project is to accurately identify bird species
from input images by developing a machine learning-based system.

To ensure the effectiveness and reliability of the developed system, a systematic
comparison analysis of multiple machine learning algorithms and deep learning
architectures will be conducted~\cite{lolge2025evolution, bilgin2022bird,
tang2025comparative}. This analysis aims to examine the performance of the algorithms.
Models such as Visual Geometry Group (VGG16), Residual Networks (ResNet models),
InceptionV3, MobileNet, and DenseNet121 will be selected~\cite{lolge2025evolution,
bilgin2022bird, tang2025comparative}.

The final outcome of this project is to build a complete citizen science platform 
that combines bird observation sharing and automated specie identification with the
help of the developed machine learning model. The project aims to contribute to biodiversity 
monitoring and the practical use of machine learning for citizen science. Platform will be used
to collect valuable ecological data that can be used by researchers

\subsection{Research Questions}

\textbf{RQ1:} How can integrating a machine learning-based bird species classification system into a citizen science platform support large-scale ecological monitoring?

\textbf{RH1:} Integrating an automated bird species classification model into a citizen science platform will improve large-scale ecological monitoring by enabling users to upload automatically labelled bird observations efficiently. This integration will reduce dependence on expert validation and increase the scalability of biodiversity data collection.

\textbf{RQ2:} How suitable are publicly available bird image datasets for training deep learning models for bird species classification?

\textbf{RH2:} Publicly available bird image datasets collected from citizen science sources such as GBIF are suitable for training bird species classification models after applying preprocessing techniques such as data cleaning, augmentation, normalization, and class balancing.